% Options for packages loaded elsewhere
\PassOptionsToPackage{unicode}{hyperref}
\PassOptionsToPackage{hyphens}{url}
\PassOptionsToPackage{dvipsnames,svgnames,x11names}{xcolor}
\PassOptionsToPackage{space}{xeCJK}
\documentclass[
]{article}
\usepackage{xcolor}
\usepackage[margin=2.5cm]{geometry}
\usepackage{amsmath,amssymb}
\setcounter{secnumdepth}{-\maxdimen} % remove section numbering
\usepackage{iftex}
\ifPDFTeX
  \usepackage[T1]{fontenc}
  \usepackage[utf8]{inputenc}
  \usepackage{textcomp} % provide euro and other symbols
\else % if luatex or xetex
  \usepackage{unicode-math} % this also loads fontspec
  \defaultfontfeatures{Scale=MatchLowercase}
  \defaultfontfeatures[\rmfamily]{Ligatures=TeX,Scale=1}
\fi
\usepackage{lmodern}
\ifPDFTeX\else
  % xetex/luatex font selection
    \setmainfont[]{DejaVu Sans}
    \setmonofont[]{DejaVu Sans Mono}
  \ifXeTeX
    \usepackage{xeCJK}
    \setCJKmainfont[]{Noto Sans CJK SC}
          \fi
  \ifLuaTeX
    \usepackage[]{luatexja-fontspec}
    \setmainjfont[]{Noto Sans CJK SC}
  \fi
\fi
% Use upquote if available, for straight quotes in verbatim environments
\IfFileExists{upquote.sty}{\usepackage{upquote}}{}
\IfFileExists{microtype.sty}{% use microtype if available
  \usepackage[]{microtype}
  \UseMicrotypeSet[protrusion]{basicmath} % disable protrusion for tt fonts
}{}
\makeatletter
\@ifundefined{KOMAClassName}{% if non-KOMA class
  \IfFileExists{parskip.sty}{%
    \usepackage{parskip}
  }{% else
    \setlength{\parindent}{0pt}
    \setlength{\parskip}{6pt plus 2pt minus 1pt}}
}{% if KOMA class
  \KOMAoptions{parskip=half}}
\makeatother
\usepackage{longtable,booktabs,array}
\usepackage{calc} % for calculating minipage widths
% Correct order of tables after \paragraph or \subparagraph
\usepackage{etoolbox}
\makeatletter
\patchcmd\longtable{\par}{\if@noskipsec\mbox{}\fi\par}{}{}
\makeatother
% Allow footnotes in longtable head/foot
\IfFileExists{footnotehyper.sty}{\usepackage{footnotehyper}}{\usepackage{footnote}}
\makesavenoteenv{longtable}
\setlength{\emergencystretch}{3em} % prevent overfull lines
\providecommand{\tightlist}{%
  \setlength{\itemsep}{0pt}\setlength{\parskip}{0pt}}
\usepackage{bookmark}
\IfFileExists{xurl.sty}{\usepackage{xurl}}{} % add URL line breaks if available
\urlstyle{same}
\hypersetup{
  colorlinks=true,
  linkcolor={Maroon},
  filecolor={Maroon},
  citecolor={Blue},
  urlcolor={Blue},
  pdfcreator={LaTeX via pandoc}}

\author{}
\date{}

\begin{document}

\section{Neural Macro Compilation in the Unified Framework of
Representation, Logic, and Intuition --- Intuition Searches
Form}\label{neural-macro-compilation-in-the-unified-framework-of-representation-logic-and-intuition-intuition-searches-form}

\section{统一框架 · 神经宏编译 ·
直觉搜索形式}\label{ux7edfux4e00ux6846ux67b6-ux795eux7ecfux5b8fux7f16ux8bd1-ux76f4ux89c9ux641cux7d22ux5f62ux5f0f}

\begin{quote}
\textbf{Draft v0.1 (English)} \textbar{} 2026-08-11 Target: ICML/COLM
(method paper) \textbar{} Patent: M1 (semantic-path compilation)
\textbf{Framework position}: Part of the \emph{Unified Framework of
Representation, Logic, and Intuition} (manifesto:
\texttt{Unified\_Framework\_Representation\_Logic\_Intuition.md}). This
paper establishes the \textbf{construction→intuition} method: compiling
the slow construction path into the fast reasoning path within one
model. Companion: P1 (form→construction theory), P0 (three-channel
equivalence cognition). Evidence: exp50/51/80/90 via run\_exp official
judge (docs/paper\_data/)
\end{quote}

\begin{center}\rule{0.5\linewidth}{0.5pt}\end{center}

\textbf{Repository:
\href{https://github.com/YuchenWang-ai/Neural-Macro-Compilation-Compiling-Slow-Construction-Paths-into-Fast-Reasoning-Paths}{YuchenWang-ai/Neural-Macro-Compilation-Compiling-Slow-Construction-Paths-into-Fast-Reasoning-Paths}}

\subsection{Abstract}\label{abstract}

A model that has learned base operations can often \emph{already
execute} a new composite operator by unfolding its definition chain ---
but along a long construction path (many steps/tokens). We propose
\emph{neural macro compilation}: compile the slow construction path into
a fast reasoning path within the \emph{same model}, using few
automatically generated samples, while (i) preserving semantics exactly
(verifiable with fallback) and (ii) never expanding the semantic
closure. On a 2-layer, \textless1 MB transformer: a 2000-digit carry
chain executes identically on construction (0.62 s), formal (0.069 s),
and intuition (0.086 s) channels (1.000); the compiled fast path
generalizes with zero decay to 2000 digits, radix 60, and 100\% digit
anonymization; step-skipping is visible as 82→5 tokens for a 20-digit
addition. Unlike distillation, teacher and student are the same model's
slow and fast paths (construction-channel self-compilation); unlike
program synthesis, no program is compiled out-of-band. Within the
Unified Framework, this paper is the \textbf{construction→intuition}
method leg: precise logical construction is compiled into flashing
intuition, which then searches forms --- closing the cycle of form →
construction → intuition → form.

\subsection{1 Introduction}\label{introduction}

\textbf{The problem.} When a model has learned base operations, a new
operator σ is already executable: its definition chain unfolds into base
operations the model can follow. But the path is long --- per-digit
iteration, long token sequences --- so inference cost scales with
construction length, not semantic size.

\textbf{Existing remedies are invasive.} Distillation needs a separate
teacher and new parameters. Few-shot fine-tuning risks disturbing
learned abilities and gives no equivalence guarantee. Program-to-neural
compilation (ICML 2024) treats the program as external input and
compiles out-of-band. None exploits that the model \emph{already} has
the semantics.

\textbf{This work.} The same model already contains two paths for σ: the
slow \emph{construction} path (definition-driven, semantically complete)
and a fast path that training can establish in weights. We compile the
former into the latter --- \emph{neural macro compilation} --- with
minimal supervision generated automatically by the slow path itself, a
verifier, and a runtime fallback. The macro adds speed, not semantics.
Relation to P1: P1 establishes \emph{that} the representation supports
systematic generalization and \emph{why} (formalized induction); this
paper establishes \emph{how} an already-executable operator is
re-compiled into a cheap path --- the two are complementary, and P2 does
not re-derive P1's generalization results.

\textbf{Contributions.} - \textbf{C1 --- Semantic-path
self-compilation.} Compile the model's own slow construction path into a
fast reasoning path with few self-generated samples, preserving
semantics and the semantic closure. - \textbf{C2 --- Zero-sample
semantic closure.} New operators are executable with zero training
(construction path); training only reduces cost. - \textbf{C3 ---
Verifiable equivalence + fallback.} Fast outputs checked against slow
outputs; failures route back to the slow path. - \textbf{C4 ---
Empirical profile.} 2000-digit carry chain correct on all channels; fast
path 7× cheaper at width 2000; zero-decay generalization; 82→5 token
step-skipping.

\subsection{2 Setting}\label{setting}

This paper builds directly on the token-native representation of P1:
concepts carry structured definitions; a \emph{definition chain} of σ is
the transitive closure of concepts σ's definition references (P1 §2).
The compilation targets are channels over this representation.

\textbf{Slow path (construction channel).} Definition-driven unfolding:
σ's definition chain expands into base operations executed step by step
(per-digit iteration, carry propagation). Semantically complete --- the
semantics is guaranteed by definition; cost scales with construction
length.

\textbf{Fast path (intuition channel).} Direct execution compiled into
weights: input → output in a short sequence, skipping intermediate
construction steps. Cost scales with output length, not construction
depth.

\textbf{Formal path (rewriting channel).} The definition-driven
rewriting engine --- shares its definition source with construction;
serves as an independent cross-check in verification.

\textbf{Goal.} For a new operator σ with expansion path T\_σ: -
Sem\_fast(σ) = Sem\_expanded(T\_σ) --- exact semantic equivalence,
verifiable; - Cost\_fast ≪ Cost\_expanded --- steps/tokens/latency
reduced; - Closure\_fast(σ) = Closure\_expanded(σ) --- no
semantic-closure expansion.

\textbf{Terminology.} The fast path is the \emph{reasoning fast path}
--- intuition: hard to obtain (compilation requires sufficient training;
config-dependent epochs), but easy to generalize by skipping steps. The
full account is in P1 (§3, §7); here we focus on the compilation method.

\subsection{3 Method: Neural Macro
Compilation}\label{method-neural-macro-compilation}

Five stages.

\textbf{S1 --- Path recognition.} Expand σ's definition chain into base
operations (E(σ) = T\_σ). If σ is expressible via the model's base
operations, it is compilable --- decided by the definition graph, not a
separate synthesizer.

\textbf{S2 --- Sample generation.} Use the slow path (construction
channel / reference evaluator) to generate (σ, x\_i, y\_i) pairs
automatically. Self-supervised labels: the samples are the model's own
semantics. Few samples suffice (tens) because the model already knows
the semantics; the samples teach only the \emph{shape} of the direct
path. The dependence of compilation quality on sample count is a
falsifiable claim, tested in §5 (currently pending).

\textbf{S3 --- Parameter update.} Update a subset of parameters
(embedding / adapter / partial / full --- graded narrow to wide) so σ
maps directly to its result in one short sequence, preserving other
abilities (restricted range, or verification after full update).

\textbf{S4 --- Semantic verification.} Verify fast-path outputs against
slow-path outputs on a held-out set by full output-sequence comparison
(the only trusted metric). Only a passing macro is installed.

\textbf{S5 --- Runtime routing.} Route σ to the fast path when its
execution-cost threshold is met; fall back to the slow path on
verification failure. The fast path is an optimization, not a
replacement.

\textbf{Key distinction from distillation.} In distillation, a (possibly
external) teacher transfers knowledge to a student. Here teacher and
student are the \emph{same model's} slow and fast paths: the
construction channel supervises its own intuition channel. No new
semantic machinery; the macro is compiled from existing competence.

\subsection{4 Experiments}\label{experiments}

All via the official run\_exp judge path (full-sequence reconstruction).
Model: exp02\_supervised\_s2 (2-layer, dim-64, \textless1 MB, judgment
1.000 on the full suite). Positioning: P1 establishes that this model
generalizes (width/radix/anonymization); this paper establishes that its
\emph{already-generalizing} capability can be re-compiled into a cheaper
path --- the experiments below measure executability, cost, and
step-skipping, not generalization itself (which P1 covers).

\subsubsection{4.1 Zero-sample semantic closure
(C2)}\label{zero-sample-semantic-closure-c2}

\begin{longtable}[]{@{}llll@{}}
\toprule\noalign{}
Input & Construction & Formal & Intuition (judgment) \\
\midrule\noalign{}
\endhead
\bottomrule\noalign{}
\endlastfoot
12+34 & ✓ & ✓ & ✓ \\
123456789+987654321 & ✓ & ✓ & ✓ \\
999\ldots9 (20 digits)+1 & ✓ & ✓ & ✓ \\
\textbf{999\ldots9 (2000 digits)+1} & ✓ & ✓ & \textbf{1.000} \\
\end{longtable}

Construction and formal channels agree exactly on all four inputs (full
output sequences); intuition (the model) scores judgment accuracy 1.000
on the same inputs, including the 2000-digit carry chain --- with zero
training for that specific operator. Semantics precede training: the
definition chain alone makes σ executable.

\subsubsection{4.2 Cost profile (C1, C4)}\label{cost-profile-c1-c4}

\begin{longtable}[]{@{}
  >{\raggedright\arraybackslash}p{(\linewidth - 8\tabcolsep) * \real{0.2000}}
  >{\raggedright\arraybackslash}p{(\linewidth - 8\tabcolsep) * \real{0.2000}}
  >{\raggedright\arraybackslash}p{(\linewidth - 8\tabcolsep) * \real{0.2000}}
  >{\raggedright\arraybackslash}p{(\linewidth - 8\tabcolsep) * \real{0.2000}}
  >{\raggedright\arraybackslash}p{(\linewidth - 8\tabcolsep) * \real{0.2000}}@{}}
\toprule\noalign{}
\begin{minipage}[b]{\linewidth}\raggedright
Width
\end{minipage} & \begin{minipage}[b]{\linewidth}\raggedright
Construction (slow)
\end{minipage} & \begin{minipage}[b]{\linewidth}\raggedright
Formal
\end{minipage} & \begin{minipage}[b]{\linewidth}\raggedright
Intuition (fast)
\end{minipage} & \begin{minipage}[b]{\linewidth}\raggedright
Speedup vs construction
\end{minipage} \\
\midrule\noalign{}
\endhead
\bottomrule\noalign{}
\endlastfoot
2 & 0.0005 s & 0.0012 s & 0.0018 s & --- \\
9 & 0.0031 s & 0.0007 s & 0.0043 s & --- \\
\textbf{2000} & \textbf{0.6206 s} & \textbf{0.0695 s} & \textbf{0.0857
s} & \textbf{7× (intuition) / 9× (formal)} \\
\end{longtable}

At 2000 digits, the fast path is 7× cheaper than construction.
Construction grows linearly with digits; the fast path is a short single
forward pass. Honest caveat: the formal engine is faster still; the
intuition channel's value is end-to-end execution in weights (no
external oracle) plus preserved generalization.

\subsubsection{4.3 Step-skipping generalization
(C4)}\label{step-skipping-generalization-c4}

At fixed acquisition (8 epochs), the compiled fast path generalizes by
skipping construction steps with \textbf{zero decay} (exp80): 2000-digit
width (10× training width), radix 60 (6×), 100\% digit anonymization,
and joint width × radix --- all at 1.000. Acquisition and step-skipping
generalization are decoupled: the fast path is hard to obtain but easy
to extrapolate.

\subsubsection{4.4 Machine form of step-skipping
(C4)}\label{machine-form-of-step-skipping-c4}

\begin{longtable}[]{@{}lll@{}}
\toprule\noalign{}
Path & Tokens & Time \\
\midrule\noalign{}
\endhead
\bottomrule\noalign{}
\endlastfoot
Construction (digit-wise unfolding) & 82 & 0.410 ms \\
Intuition (atomic num token) & 5 & 0.186 ms \\
\end{longtable}

20-digit addition: the fast path treats the whole numeral as an atomic
token (5 tokens) versus 82-token digit-wise unfolding --- a 16× token
reduction, 2.2× faster. \textbf{Honest caveat}: the atomic token is
outside the model vocab (pad id 0); the measurement is the
sequence-length effect on forward cost --- a lower bound, not a claim of
learned atomic-token semantics. If the model truly learned the atomic
token, the fast path would be an O(1) decision with a larger advantage.

\subsection{5 Ablations}\label{ablations}

\begin{longtable}[]{@{}
  >{\raggedright\arraybackslash}p{(\linewidth - 4\tabcolsep) * \real{0.3333}}
  >{\raggedright\arraybackslash}p{(\linewidth - 4\tabcolsep) * \real{0.3333}}
  >{\raggedright\arraybackslash}p{(\linewidth - 4\tabcolsep) * \real{0.3333}}@{}}
\toprule\noalign{}
\begin{minipage}[b]{\linewidth}\raggedright
Ablation
\end{minipage} & \begin{minipage}[b]{\linewidth}\raggedright
Question
\end{minipage} & \begin{minipage}[b]{\linewidth}\raggedright
Status
\end{minipage} \\
\midrule\noalign{}
\endhead
\bottomrule\noalign{}
\endlastfoot
Sample count (10/20/50) & how few instances suffice for compilation &
pending \\
Update scope (embeddings / adapter / partial / full) & cost--degradation
trade-off; original-capability protection & pending \\
No verifier / no fallback & error propagation without the safety net &
pending \\
Closure invariance (EXP-52) & fast-path executable set = slow-path
executable set & pending \\
\end{longtable}

The three-channel consistency experiment (exp50) already provides the
core verification evidence: construction = formal exactly (2/9/20/2000
digits) and intuition 1.000, including the 2000-digit carry chain.
Closure invariance (EXP-52) tests whether the fast path executes exactly
the inputs the slow path could execute --- the semantic-closure
guarantee of Goal (iii).

\subsection{6 Related Work}\label{related-work}

\begin{itemize}
\tightlist
\item
  \textbf{Program-to-neural compilation} (Learning to Compile Programs
  to Neural Networks, ICML 2024): builds a neural surrogate for an
  \emph{external} program from behavioral data; requires a program
  specification and large data. Difference: here the ``program'' is the
  model's own construction path; the teacher is the model itself; no
  program analysis. Related survey perspective in P1 §9.
\item
  \textbf{Knowledge distillation} (Hinton et al., 2015):
  teacher→student, typically a separate/larger teacher. Difference: same
  model, two channels; no external teacher; equivalence verified with
  fallback.
\item
  \textbf{Inference compilation / shortcut distillation} (diffusion,
  2025): few-sample shortcut distillation compresses long correct paths.
  Similar surface; difference: semantic equivalence with a verifier +
  fallback, no semantic-closure expansion.
\item
  \textbf{Few-shot fine-tuning}: general; this work is the special case
  where the capability already exists in the slow path, so few instances
  \emph{compile} rather than \emph{learn}.
\item
  \textbf{Continuous / continual learning}: adding a capability without
  forgetting old ones. This work's update-scope restriction (S3) and
  verification (S4) are a forgetting-mitigation mechanism; the ablation
  (§5) measures the trade-off.
\item
  \textbf{Mechanistic / hand-compiled transformers} (e.g.~a
  hand-compiled WASM interpreter): hand-assigned weights interpret a
  language. Ours is \emph{learned} compilation --- gradient training
  produces the macro, no hand-set weights.
\item
  \textbf{Continual learning / catastrophic forgetting}: fine-tuning
  risks forgetting; our restricted-update + verification + fallback is a
  targeted response, and the update-scope ablation (§5) tests it
  explicitly.
\item
  \textbf{Length generalization literature} (positional encoding /
  recurrence engineering): addresses long-input inference under fixed
  representations; this work's fast path removes construction-length
  dependence at the channel level rather than the representation level.
\end{itemize}

\subsection{7 Limitations}\label{limitations}

\begin{enumerate}
\def\labelenumi{\arabic{enumi}.}
\tightlist
\item
  \textbf{Compilable set.} Only operators expressible as compositions of
  already-supported operations can be compiled (S1 precondition). New
  semantics are out of scope: the macro adds speed, not meaning.
\item
  \textbf{Slow-path correctness.} Compilation quality inherits the
  construction channel's correctness; a reference-evaluator bug
  propagates into supervision. Verification mitigates but does not
  create correctness.
\item
  \textbf{Verification coverage.} Equivalence is checked on held-out
  samples, not proven; a full guarantee would require formal
  verification of the compiled weights (future, Lean-based).
\item
  \textbf{Update-range tradeoff.} Wide updates risk disturbing learned
  abilities; narrow updates may limit the fast path's expressiveness.
  The empirical frontier is not yet mapped (ablation pending).
\end{enumerate}

\subsection{8 Conclusion}\label{conclusion}

A model that already executes a composite operator along its
construction path has everything needed to execute it fast: a small
number of samples compiles the construction channel into a semantically
equivalent reasoning fast path within the same model --- verified, with
fallback, and without expanding the semantic closure. On a 2-layer
\textless1 MB transformer: a 2000-digit carry chain is correct on all
channels, the fast path is 7× cheaper at width 2000, step-skipping
generalizes with zero decay, and the macro's machine form is visible as
82→5 tokens. Neural macro compilation is the weight-level analogue of
source-level macros: the semantics are the model's own; only the speed
is added.

\subsection{Appendix}\label{appendix}

\begin{itemize}
\tightlist
\item
  A. Full experiment tables: docs/paper\_data/results.csv.
\item
  B. Path-recognition examples: definition-chain expansions for
  addition/root/tetration.
\item
  C. Cost protocol: construction (token\_iterator), formal (engine),
  intuition (model forward); wall-clock per 1000 samples.
\item
  D. Patent M1 independent-claim skeleton (semantic-path compilation).
\end{itemize}

\subsection{Cross-References within the Unified
Framework}\label{cross-references-within-the-unified-framework}

This paper is one leg of the \emph{Unified Framework of Representation,
Logic, and Intuition} (manifesto:
\texttt{Unified\_Framework\_Representation\_Logic\_Intuition.md}). The
three legs co-adapt:

\begin{longtable}[]{@{}
  >{\raggedright\arraybackslash}p{(\linewidth - 4\tabcolsep) * \real{0.3333}}
  >{\raggedright\arraybackslash}p{(\linewidth - 4\tabcolsep) * \real{0.3333}}
  >{\raggedright\arraybackslash}p{(\linewidth - 4\tabcolsep) * \real{0.3333}}@{}}
\toprule\noalign{}
\begin{minipage}[b]{\linewidth}\raggedright
Leg
\end{minipage} & \begin{minipage}[b]{\linewidth}\raggedright
Paper
\end{minipage} & \begin{minipage}[b]{\linewidth}\raggedright
Claim
\end{minipage} \\
\midrule\noalign{}
\endhead
\bottomrule\noalign{}
\endlastfoot
form → construction & P1 (this paper) & correct syntax → extreme
training speed; correct construction → extreme generalization \\
construction → intuition & P0 (three-channel equivalence) & construction
= formal = intuition; intuition is compiled construction \\
construction → intuition (method) & P2 (neural macro compilation) &
compile the slow construction path into the fast reasoning path \\
\end{longtable}

Cyclic closure: form directs construction (P1), construction earns
intuition (P0/P2), intuition searches form (the fast path extrapolates
and re-expresses in the unified form). The unified conclusion --- form
directs logical construction, construction earns flashing intuition,
intuition exhaustively searches forms --- is stated in the manifesto and
in each paper's Unified Conclusion.

\subsection{References}\label{references}

{[}In preparation; verified full records pending. Core entries:{]}

\begin{itemize}
\tightlist
\item
  Hinton, G., Vinyals, O., Dean, J. 2015. Distilling the Knowledge in a
  Neural Network. arXiv:1503.02531.
\item
  ICML 2024. Learning to Compile Programs to Neural Networks (full
  citation pending verification).
\item
  P1 companion paper: Generalization as Formalized Induction (2026).
\item
  Diffusion shortcut distillation (2025; citation pending).
\item
  Barendregt, H. 1984. The Lambda Calculus (for definition-chain and λ
  foundations).
\end{itemize}

\begin{center}\rule{0.5\linewidth}{0.5pt}\end{center}

\subsection{Unified Conclusion (added
2026-08-11)}\label{unified-conclusion-added-2026-08-11}

Once, under a unified formal attention-syntactic structure, a neural
network ultimately, inevitably, and necessarily clearly, intuitively,
and interpretably observes the boundary between intuition and
construction, we can with full confidence use the unified
attention-formal language to direct logical construction, use precise
synthetic-CoT logical construction to earn \emph{flashing intuition},
and use distilled intuitive context to exhaustively search attention
forms. Form, construction, and intuition co-adapt within the unified
framework, forming a feedback iteration pipeline of teaching, training,
and reasoning in which they are no longer distinct --- perhaps one of
the faster shortcuts among the many paths toward ASI.

\end{document}
